\subsection{Teacher Knowledge, Qualifications, and the Teaching of Algorithms}

Research on computer science education increasingly emphasizes that the quality of algorithm instruction in secondary education is closely tied to
  teacher preparation, pedagogical content knowledge, and professional development.
Due to the inherent difficulties of teaching algorithms, teachers need more than general classroom skill or basic familiarity with programming
  \cite{yadav2016computational}.
Doukakis et al.
found that Greek secondary computer science teachers who taught algorithms and programming reported strong content knowledge and technological knowledge, but lower confidence in pedagogical content knowledge and technological content knowledge \cite{doukakis2013measuring}.
This suggests that even formally prepared computer science teachers may feel less certain about transforming algorithmic content into teachable
  representations, activities, and explanations for students.
Giannakos et al.
reinforces this, who examined a national sample of Greek computer science teachers through the TPACK framework \cite{giannakos2014examining}.
Their findings showed that teachers rated their content knowledge highly, but reported weaker confidence in pedagogical content knowledge, especially
  in relation to teaching algorithms and integrating educational software.
Notably, 43\% of the teachers in the study indicated a need for additional training specifically in how to teach algorithms
  \cite{giannakos2014examining}.
Teachers may understand algorithms themselves, but still need support in selecting examples, anticipating misconceptions, scaffolding
  problem-solving.

Leyzberg and Moretti’s work focuses on professional development for secondary CS teachers \cite{leyzberg2017teaching}.
They note that many U.S.
computer science teachers do not have formal CS degrees, and that introductory professional development often focuses on helping teachers deliver first-exposure CS education rather than developing advanced content knowledge and understanding.
Their workshop model addressed topics such as recursion, algorithm and data structure analysis, theoretical computer science, and abstraction, all of
  which are closely related to teaching algorithms at a deeper level.
Their findings suggest that teachers value rigorous content-focused professional development, especially when it helps them support advanced students
  and move beyond shallow or purely procedural instruction.

Literature on the subject show’s correlation through compounding factors in the rapid expansion of secondary computer science education.
Zhou et al.
noted that increasing enrollment in secondary CS courses has created a shortage of well-prepared teachers capable of teaching computational concepts effectively \cite{zhou2020high}.
Their research on the subject emphasizes that computer science content is inherently difficult and that teaching equitable, accessible CS lessons
  requires both substantial content knowledge and strong pedagogical confidence.
Algorithms require persistence, iterative problem solving, and comfort with ambiguity, all of which can be difficult for non-specialized teachers to
  facilitate successfully \cite{zhou2020high}.

Research completed by Nijenhuis-Voogt et al.
studied teachers’ pedagogical content knowledge for teaching algorithms in upper secondary education \cite{nijenhuis2023teaching}.
Their findings show that teachers vary in whether they frame algorithms primarily as objects for “thinking” or as objects for both “thinking and
  making,” meaning that some teachers emphasize analysis and reasoning while others also emphasize implementation in code.
The study also found that teachers need topic-specific knowledge for responding to student differences, including how to scaffold students who
  struggle to begin solving algorithmic problems, generalize from examples, or analyze efficiency \cite{nijenhuis2023teaching}.
This makes PCK especially important for algorithm instruction, as algorithmic thinking involves both conceptual understanding and strategic
  problem-solving.

Zhou et al.
extend this discussion by focusing on teacher self-efficacy in computer science teaching.
Their study found that a hybrid professional development program significantly increased teachers’ self-efficacy in both CS content knowledge and
  pedagogical content knowledge \cite{zhou2020high}.
Demonstrating that self-efficacy can influence whether teachers persist through difficult instructional challenges, experiment with student-centered
  strategies, and whether they are able to support students through the iterative problem-solving process that algorithms require.
Their work suggests that professional development should give teachers opportunities for mastery experiences, peer collaboration, and practice with
  CS-specific instructional strategies rather than relying only on general pedagogy \cite{zhou2020high}.

Related to this is a disagreement over the value of prior programming experience.
Some teachers in Nijenhuis-Voogt et al.
viewed programming experience positively because it helped students think more systematically and understand structures such as iteration and selection \cite{nijenhuis2023teaching}.
Other teachers believed prior programming experience could hinder algorithmic thinking because students would rush into coding solutions instead of
  analyzing the problem carefully first \cite{futschek2006algorithmic}\cite{bell2014presenting}\cite{basu2025using}\cite{Oshanova2019}.

Research also disagrees with how teacher qualifications are conceptualized.
Some papers place stronger emphasis on content knowledge, while others emphasize pedagogical content knowledge (PCK).
Leyzberg and Moretti argue that many teachers need deeper exposure to subjects such as recursion, abstraction, and algorithm analysis
  \cite{leyzberg2017teaching}.
In contrast, the TPACK-oriented studies by Doukakis et al.
and Giannakos et al. report that educators require more pedagogical knowledge rather than PCK \cite{doukakis2013measuring}\cite{giannakos2014examining}.
Zhou et al.
supports this through their stance on teacher self-efficacy and confidence-building through collaborative professional development experiences \cite{zhou2020high}.
They argue that teachers’ beliefs about their abilities strongly influence classroom practice.
Meanwhile, other studies focus more heavily on concrete knowledge domains such as algorithm analysis, instructional scaffolding, or pedagogical
  techniques.
